\chapter{Parathyroidectomy}

\section*{Introduction \& Clinical Indications}
Parathyroidectomy is a surgical intervention aimed at the removal of one or more parathyroid glands, which are small endocrine glands located in the neck, typically posterior to the thyroid gland. This procedure is primarily indicated for the treatment of primary hyperparathyroidism, a condition characterized by excessive secretion of parathyroid hormone (PTH), leading to hypercalcemia and a spectrum of clinical manifestations. The surgical goal is to excise the hyperfunctioning gland(s), thereby restoring normal calcium levels and alleviating symptoms such as bone pain, kidney stones, and neurocognitive disturbances. Parathyroidectomy is recognized as a definitive treatment for hyperparathyroidism, with a high success rate in normalizing serum calcium levels and improving patient quality of life.

\begin{itemize}
  \item Parathyroid gland removal
  \item Parathyroid surgery
  \item Minimally invasive parathyroidectomy (MIP)
  \item Focused parathyroidectomy
  \item Bilateral neck exploration for hyperparathyroidism
  \item Parathyroid adenoma excision
\end{itemize}

The primary surgical principle of parathyroidectomy is the accurate identification and complete excision of the gland(s) responsible for excessive secretion of parathyroid hormone (PTH). In the context of primary hyperparathyroidism, this is most often an autonomously functioning parathyroid adenoma, although multiglandular hyperplasia or parathyroid carcinoma may also be involved. The rationale for surgical intervention is based on the adverse effects of chronic hypercalcemia and elevated PTH levels, which can lead to significant complications, including osteoporosis, renal stones, and neuropsychiatric symptoms.

By excising the hyperfunctioning gland(s), the normal feedback mechanism between serum calcium and PTH is restored, allowing for the normalization of calcium levels and alleviation of associated symptoms. The technical goals of the procedure include achieving a biochemical cure of hyperparathyroidism, minimizing surgical morbidity, and preserving parathyroid function to prevent permanent hypoparathyroidism. Intraoperative parathyroid hormone (IOPTH) monitoring is a valuable adjunct, providing real-time biochemical confirmation of successful gland removal by demonstrating a significant decrease in PTH levels during the procedure.

The history of parathyroidectomy reflects significant advancements in surgical techniques and understanding of endocrine disorders. The first successful parathyroidectomy was performed by Felix Mandl in 1925, who excised a parathyroid adenoma from the mediastinum of a patient suffering from severe osteitis fibrosa cystica. This pioneering surgery established parathyroidectomy as a viable treatment for hyperparathyroidism.

In the following decades, surgical approaches evolved from extensive bilateral neck explorations to more targeted techniques. The advent of reliable imaging modalities in the 1970s and 1980s, such as technetium-99m sestamibi scans and high-resolution ultrasound, greatly improved the localization of parathyroid glands. The introduction of intraoperative PTH monitoring in the 1990s further refined the procedure, allowing for confirmation of successful gland removal and reducing the need for extensive exploration. Today, minimally invasive techniques are often employed, making parathyroidectomy a safe and effective procedure with high success rates.

Parathyroidectomy is primarily indicated for patients diagnosed with primary hyperparathyroidism who meet specific criteria, as well as certain cases of secondary and tertiary hyperparathyroidism.

\begin{itemize}
  \item Symptomatic primary hyperparathyroidism, including bone pain, fatigue, generalized weakness, depression, or neurocognitive dysfunction.
  \item Significant hypercalcemia, typically defined as serum calcium levels greater than 1.0 mg/dL above the upper limit of normal, even in asymptomatic patients.
  \item Presence of kidney stones (nephrolithiasis) or nephrocalcinosis attributed to hyperparathyroidism.
  \item Reduced bone mineral density, such as osteoporosis (T-score less than -2.5) or osteopenia, particularly at sites affected by PTH (e.g., distal radius).
  \item Age less than 50 years, as younger patients are more likely to develop complications over their lifetime.
  \item Creatinine clearance reduced by more than 30\% compared to age-matched controls, indicating renal impairment.
  \item Urinary calcium excretion greater than 400 mg/24 hours, increasing the risk of kidney stone formation.
  \item Development of a parathyroid crisis, a rare but severe manifestation of hyperparathyroidism with extreme hypercalcemia requiring urgent intervention.
  \item Certain cases of secondary hyperparathyroidism in patients with chronic kidney disease who fail medical management and develop severe bone disease, intractable pruritus, or calciphylaxis.
  \item Tertiary hyperparathyroidism, where parathyroid glands become autonomous after prolonged secondary hyperparathyroidism, often following successful renal transplantation.
\end{itemize}

Parathyroidectomy is considered the primary and definitive curative treatment for primary hyperparathyroidism. For the majority of patients meeting surgical criteria, it is the first-line intervention with exceptionally high success rates in normalizing calcium and PTH levels, leading to resolution of symptoms and prevention of long-term complications. Medical management, primarily with calcimimetics (e.g., cinacalcet) or bisphosphonates, is typically reserved for patients who are not surgical candidates due to significant comorbidities, those who decline surgery, or those with mild, asymptomatic disease who are closely monitored.

In the context of secondary hyperparathyroidism, which arises from chronic kidney disease, parathyroidectomy serves as a secondary or salvage procedure. It is performed when medical therapies (e.g., vitamin D analogs, phosphate binders, calcimimetics) fail to control severe hyperparathyroidism, leading to intractable symptoms, severe bone disease, or calciphylaxis. For tertiary hyperparathyroidism, which often develops after successful renal transplantation when previously hyperplastic glands become autonomous, parathyroidectomy is again a primary curative intervention for the now independent parathyroid glands.

\section*{Relevant Surgical Anatomy}
The parathyroid glands are typically four small, ovoid structures, each weighing approximately 30-50 mg, located in the neck. They are usually arranged in two pairs: superior and inferior. The superior parathyroid glands are generally more consistent in their anatomical position, located posterior to the recurrent laryngeal nerve and at the level of the cricoid cartilage, often within 1 cm of the intersection of the recurrent laryngeal nerve and the inferior thyroid artery. The inferior parathyroid glands exhibit greater variability in their location, commonly found near the inferior pole of the thyroid gland, but can also descend with the thymus into the anterior mediastinum or be located within the thyroid parenchyma, carotid sheath, or even retroesophageally. 

The blood supply to the parathyroid glands is primarily derived from the inferior thyroid arteries, branches of the thyrocervical trunk. Occasionally, the superior parathyroid glands may receive additional blood supply from the superior thyroid arteries. Preservation of the vascular supply to any remaining parathyroid tissue is crucial to prevent postoperative hypoparathyroidism. The recurrent laryngeal nerves (RLN), which innervate the intrinsic muscles of the larynx, run in close proximity to the parathyroid glands, typically in the tracheoesophageal groove. The external branch of the superior laryngeal nerve (SLN), which innervates the cricothyroid muscle, also lies near the superior pole of the thyroid and is at risk during surgery. Lymphatic drainage generally follows the thyroid lymphatics. Key surgical landmarks include the inferior pole of the thyroid, Berry's ligament, and the tracheoesophageal groove.

\section*{Preoperative Workup \& Preparation}
Thorough preoperative preparation is crucial for a safe and successful parathyroidectomy, often involving a multidisciplinary approach to optimize patient outcomes. The primary goals are to confirm the diagnosis of hyperparathyroidism, accurately localize the abnormal gland(s), and assess the patient's overall medical fitness for surgery.

\begin{itemize}
  \item Diagnostic Confirmation: Repeated measurements of elevated serum calcium and parathyroid hormone (PTH) levels are essential to confirm the diagnosis of primary hyperparathyroidism.
  \item Localization Studies: Preoperative imaging is critical to guide the surgeon, commonly including Technetium-99m sestamibi scan (nuclear medicine scan), high-resolution neck ultrasound, and often four-dimensional computed tomography (4D CT) of the neck and mediastinum. Magnetic resonance imaging (MRI) may be used in select complex cases or for ectopic glands.
  \item Biochemical Assessment: Comprehensive blood work includes calcium, albumin, PTH, vitamin D (25-hydroxyvitamin D), creatinine, and electrolytes. Vitamin D deficiency should ideally be corrected preoperatively to minimize the risk of postoperative hypocalcemia.
  \item Anesthesia Consultation: Evaluation by an anesthesiologist to assess overall health, identify potential risks, and plan for general anesthesia, especially for patients with significant comorbidities.
  \item Cardiac and Pulmonary Clearance: For patients with significant cardiovascular or pulmonary comorbidities, cardiac stress testing, echocardiography, or pulmonary function tests may be required to ensure surgical fitness.
  \item Medication Review: Adjustment or temporary discontinuation of certain medications, such as blood thinners (e.g., warfarin, novel oral anticoagulants) or medications affecting calcium metabolism (e.g., lithium), is necessary.
  \item NPO Status: Patients are instructed to be NPO (nil per os) for a specified period (typically 6-8 hours) before surgery to prevent aspiration during anesthesia.
  \item Antibiotic Prophylaxis: Administration of prophylactic antibiotics before incision is standard practice to reduce the risk of surgical site infection.
  \item Informed Consent: A detailed discussion with the patient regarding the procedure, potential risks, benefits, and alternative treatments, followed by obtaining informed consent, is mandatory.
\end{itemize}

While parathyroidectomy is generally safe and highly effective, certain conditions may contraindicate or necessitate extreme caution during the procedure.

\textbf{Absolute Contraindications:}
\begin{itemize}
  \item Uncontrolled severe medical comorbidities that make general anesthesia or surgery excessively risky, such as recent myocardial infarction, unstable angina, severe uncontrolled heart failure, or severe respiratory failure.
  \item Inability to tolerate general anesthesia due to severe allergic reactions or other physiological limitations.
  \item Active, uncontrolled infection in the surgical field or systemic sepsis.
\end{itemize}

\textbf{Relative Contraindications and Precautions:}
\begin{itemize}
  \item Poor Localization: Inability to definitively localize the abnormal gland(s) preoperatively can make a focused approach challenging and may necessitate a more extensive bilateral neck exploration, increasing operative time and potential morbidity.
  \item Prior Neck Surgery or Radiation: Previous neck operations (e.g., thyroidectomy, carotid endarterectomy) or radiation therapy can lead to significant scarring and anatomical distortion, increasing the technical difficulty of dissection and the risk of complications, particularly recurrent laryngeal nerve injury.
  \item Recurrent Laryngeal Nerve Injury Risk: Meticulous dissection and, in some cases, intraoperative nerve monitoring are crucial to protect the recurrent laryngeal nerves, which are vital for vocal cord function. Injury can lead to temporary or permanent hoarseness or vocal cord paralysis.
  \item Superior Laryngeal Nerve Injury Risk: While less common, injury to the external branch of the superior laryngeal nerve can impair vocal pitch and projection, which is particularly relevant for professional singers or speakers.
  \item Ectopic Glands: Parathyroid glands can be located in ectopic positions (e.g., mediastinum, thymus, carotid sheath, intrathyroidal), making their identification and removal more challenging and potentially requiring specialized approaches or even sternotomy for mediastinal glands.
  \item Parathyroid Carcinoma: Suspected or confirmed parathyroid carcinoma requires a more aggressive surgical approach, including en bloc resection with surrounding tissues and potentially ipsilateral thyroid lobectomy, due to its invasive nature.
  \item Severe Hypercalcemia: Patients with very high calcium levels may require aggressive medical management (e.g., intravenous fluids, bisphosphonates, calcitonin) to lower calcium before surgery to reduce perioperative risks such as cardiac arrhythmias or renal failure.
\end{itemize}

\section*{Surgical Technique \& Procedure}
Parathyroidectomy is performed under general anesthesia, with the patient typically positioned supine with the neck extended by placing a shoulder roll, and the head slightly turned to the side opposite the planned incision for unilateral approaches. Neuromonitoring of the recurrent laryngeal nerve may be utilized to aid in its identification and protection.

The specific surgical approach varies based on preoperative localization studies and surgeon preference. For a focused parathyroidectomy, a small, typically 2-3 cm transverse incision is made in a natural skin crease on the side of the neck where the abnormal gland was localized. For bilateral neck exploration, a slightly longer central transverse incision is made. The platysma muscle is then divided, and the strap muscles (sternohyoid, sternothyroid, omohyoid) are carefully separated or divided to expose the thyroid gland. The surgeon meticulously dissects around the thyroid to identify the parathyroid glands, which are typically small, yellowish-brown, and ovoid, often embedded in fat. Preoperative imaging studies are crucial in guiding the search for the abnormal gland.

Once the suspected abnormal gland (adenoma) is identified, intraoperative parathyroid hormone (IOPTH) monitoring is initiated. Blood samples are drawn before incision (baseline), after gland exposure, and at 5-10 minute intervals after removal of the suspected abnormal gland. A significant drop (typically >50\%) from the highest pre-excision PTH level confirms successful removal of hyperfunctioning tissue. The abnormal gland is then carefully dissected free from surrounding tissues, ensuring meticulous hemostasis and, critically, preserving the blood supply to any remaining normal parathyroid glands. Care is taken to avoid injury to the recurrent laryngeal nerve. In cases of multiglandular hyperplasia, subtotal parathyroidectomy (removal of 3.5 glands) or total parathyroidectomy with autotransplantation of a small portion of parathyroid tissue into a forearm muscle may be performed to prevent permanent hypoparathyroidism. After successful gland removal and confirmation of PTH drop, meticulous hemostasis is achieved, the wound is irrigated, and the strap muscles are reapproximated. The platysma and skin are closed in layers, often with absorbable sutures. Drains are rarely needed for routine cases. The entire procedure typically takes 1 to 2 hours for a focused approach, and slightly longer for bilateral exploration.

\section*{Complications \& Risks}
As with any surgical procedure, parathyroidectomy carries potential risks, which can be broadly categorized into general surgical complications and procedure-specific risks.

\textbf{General Surgical Risks:}
\begin{itemize}
  \item Bleeding: Hematoma formation in the neck, which can rarely be severe enough to cause airway compromise requiring urgent re-exploration.
  \item Infection: Surgical site infection, though uncommon in clean neck surgery, can occur.
  \item Anesthesia Complications: Reactions to anesthetic agents, respiratory or cardiac complications, nausea, vomiting, or aspiration.
  \item Pain: Postoperative pain at the incision site, usually mild to moderate and manageable with oral analgesics.
  \item Scarring: Formation of a visible scar on the neck, which typically fades over time.
\end{itemize}

\textbf{Procedure-Specific Risks:}
\begin{itemize}
  \item Recurrent Laryngeal Nerve Injury: Temporary or permanent hoarseness or vocal cord paralysis due to injury to the recurrent laryngeal nerve, which runs close to the parathyroid glands. This can range from mild voice changes to significant breathing difficulties.
  \item Hypocalcemia (Low Calcium):
    \begin{itemize}
      \item Transient Hypocalcemia: Common in the immediate postoperative period as the remaining parathyroid glands adjust to the sudden drop in PTH.
      \item Permanent Hypoparathyroidism: Occurs if all parathyroid glands are inadvertently removed or devascularized, leading to lifelong dependence on calcium and vitamin D supplementation.
      \item Hungry Bone Syndrome: A severe form of hypocalcemia seen in patients with significant preoperative bone disease, where calcium rapidly shifts into the bones after PTH removal, requiring aggressive calcium and vitamin D repletion.
    \end{itemize}
  \item Persistent Hyperparathyroidism: Failure to cure hyperparathyroidism if the abnormal gland is not found or if additional hyperfunctioning glands are missed during the initial surgery. This necessitates further investigation and potentially reoperation.
  \item Recurrent Hyperparathyroidism: Development of hyperparathyroidism months or years after initial successful surgery, often due to growth of a previously normal gland or a missed adenoma that subsequently becomes hyperfunctional.
  \item Superior Laryngeal Nerve Injury: Rare injury to the external branch, leading to subtle voice changes, particularly affecting vocal pitch and projection.
  \item Seroma: Collection of fluid under the skin flap, which may require aspiration.
  \item Damage to Adjacent Structures: Rare injury to the thyroid gland, trachea, or esophagus during dissection.
\end{itemize}

\section*{Postoperative Care \& Recovery}
Immediately after parathyroidectomy, patients are transferred to the Post-Anesthesia Care Unit (PACU) for close monitoring as they recover from general anesthesia. Vital signs, pain levels, and the neck for swelling or hematoma are continuously assessed. Blood calcium levels are typically checked within a few hours post-op and then regularly over the next 24-48 hours to monitor for hypocalcemia, which is a common transient phenomenon. Patients may experience mild to moderate neck pain, which is usually managed effectively with oral analgesics, and some transient hoarseness or difficulty swallowing, which generally resolves within a few days. Most patients are able to resume oral intake of liquids and soft foods within a few hours of surgery.

For the first 1-2 weeks following discharge, patients are advised to limit strenuous activities, heavy lifting, and excessive neck movement to promote optimal wound healing. The incision site should be kept clean and dry, and specific instructions for wound care will be provided. Calcium and vitamin D supplementation may be prescribed, especially if there is a risk of hungry bone syndrome or if preoperative vitamin D deficiency was present. Most patients undergoing minimally invasive procedures are discharged within 24 hours, often on the same day. Long-term recovery involves ongoing monitoring of calcium and PTH levels, and addressing any residual symptoms or complications. Full recovery and complete resolution of symptoms, particularly bone-related issues, can take several weeks to months.

During parathyroidectomy, the surgeon meticulously documents the location, size, and appearance of all identified parathyroid glands, noting any abnormalities. A hyperfunctioning parathyroid adenoma typically appears as an enlarged, often darker, and softer gland compared to normal parathyroid tissue, which is usually small, flat, and yellowish-brown. In cases of hyperplasia, multiple glands may appear enlarged. Intraoperative frozen section analysis is often performed on excised tissue to confirm that it is indeed parathyroid tissue and to distinguish it from thyroid tissue, lymph nodes, or fat. This helps guide the surgeon in ensuring the correct tissue has been removed.

The definitive diagnosis is made by the pathologist on the resected tissue. The pathology report will confirm the nature of the excised gland(s), most commonly a parathyroid adenoma, characterized by a uniform proliferation of chief cells with minimal fat, often surrounded by a thin capsule. Parathyroid hyperplasia involves enlargement of multiple glands, while parathyroid carcinoma is a rare diagnosis characterized by features such as capsular invasion, vascular invasion, and mitotic activity, which necessitate a more aggressive surgical approach. The pathology report is crucial for confirming the cause of hyperparathyroidism and guiding long-term management.

A normal and successful postoperative course following parathyroidectomy is characterized by the rapid resolution of hyperparathyroidism and a gradual improvement in associated symptoms. Within 24-48 hours, serum calcium levels are expected to return to the normal range, often accompanied by a significant drop in parathyroid hormone (PTH) levels. Patients may experience mild neck pain and some transient hoarseness or difficulty swallowing, which typically resolve within a few days to a week. Transient hypocalcemia is a common and expected phenomenon as the remaining parathyroid glands adjust to the sudden removal of the hyperfunctioning tissue, and this is usually managed with oral calcium and vitamin D supplementation. Over the ensuing weeks to months, patients can expect a gradual improvement or complete resolution of symptoms such as fatigue, bone pain, muscle weakness, and neurocognitive issues. The surgical incision typically heals well, leaving a fine scar that fades over time, allowing for a full return to normal activities.

Postoperative care and long-term follow-up are essential to ensure sustained cure, monitor for potential complications, and address any residual issues.

\begin{itemize}
  \item Initial Postoperative Visit: Typically scheduled within 1-2 weeks after surgery for a wound check, removal of non-absorbable sutures or staples, and initial assessment of symptoms and calcium levels.
  \item Biochemical Monitoring: Serial measurements of serum calcium and PTH levels are crucial, often at 1 month, 3 months, 6 months, and then annually for several years to confirm sustained normocalcemia and rule out recurrence. Vitamin D levels may also be monitored, and supplementation continued if deficiency persists.
  \item Bone Mineral Density (BMD) Testing: A repeat DEXA scan is often recommended 1 year after successful parathyroidectomy to assess for improvement in bone density, especially if osteoporosis was present preoperatively.
  \item Renal Function Assessment: Monitoring of kidney function (creatinine, eGFR) and urine calcium excretion is important, particularly if kidney stones or renal impairment were present preoperatively.
  \item Activity Resumption: Gradual return to normal activities is encouraged, with specific guidance on when to resume strenuous exercise or heavy lifting, typically after 2-4 weeks.
  \item Warning Signs: Patients are educated on symptoms of hypocalcemia (e.g., tingling around the mouth or in fingers, muscle cramps, spasms, numbness) and instructed to contact their physician immediately if these occur. They should also report any persistent hoarseness, difficulty breathing, or signs of wound infection (redness, swelling, discharge, fever).
\end{itemize}

Long-term follow-up with an endocrinologist or the operating surgeon is crucial to ensure the durability of the cure and manage any long-term sequelae.

\section*{Outcomes \& Prognosis}
Primary hyperparathyroidism (PHPT), the most common indication for parathyroidectomy, has an estimated incidence of 21 per 100,000 person-years, though this varies significantly by age and sex. It is notably more prevalent in women, with a female-to-male ratio of approximately 3:1 to 4:1. The incidence increases with age, peaking in the 5th to 7th decades of life, with a mean age at diagnosis typically in the 60s. The vast majority of cases (80-85\%) are caused by a single parathyroid adenoma, while 10-15\% are due to multiglandular hyperplasia, and less than 1\% are caused by parathyroid carcinoma. Familial forms, such as Multiple Endocrine Neoplasia type 1 (MEN1) or type 2A (MEN2A), account for a small percentage of cases. The widespread use of automated biochemical screening panels has led to an increase in the diagnosis of asymptomatic PHPT, which now constitutes a significant proportion of cases. Mortality directly attributable to parathyroidectomy is exceedingly low, typically less than 0.1\%, reflecting the safety of modern surgical techniques. Morbidity, primarily related to transient hypocalcemia or recurrent laryngeal nerve injury, is also low, generally ranging from 1-5\%.

Parathyroidectomy offers an excellent prognosis for patients with primary hyperparathyroidism, with high rates of biochemical cure and significant improvement in symptoms. The success rate for initial parathyroidectomy, defined as normalization of serum calcium and PTH levels, is typically 95-98\% when performed by experienced surgeons, particularly with the aid of accurate preoperative localization and intraoperative PTH monitoring. For patients with a single adenoma, the cure rate is even higher. Symptomatic improvement, including reduction in bone pain, fatigue, and neurocognitive symptoms, is observed in a majority of patients, though the timeline for resolution can vary. Bone mineral density often improves significantly within 1-2 years post-surgery, reducing fracture risk. The risk of persistent hyperparathyroidism (failure to cure) is low, usually 1-3\%, and is often due to missed adenomas or unrecognized multiglandular disease. Recurrent hyperparathyroidism, defined as hypercalcemia occurring at least 6 months after successful surgery, is also uncommon, affecting 2-5\% of patients over 10-15 years, and may be due to a new adenoma or growth of previously hyperplastic tissue. Prognostic factors for successful outcome include accurate preoperative localization, surgeon experience, and the absence of multiglandular disease or carcinoma. Overall, parathyroidectomy is a highly effective and durable treatment for hyperparathyroidism, significantly improving patient quality of life and reducing long-term complications.

\section*{References}
\begin{enumerate}
  \item MedlinePlus. Parathyroidectomy. U.S. National Library of Medicine; 2024. Available from: \url{https://medlineplus.gov/ency/article/002937.htm}
  \item Doherty GM, Moley JF. Parathyroidectomy. In: Townsend CM Jr, Beauchamp RD, Evers BM, Mattox KL, editors. Sabiston Textbook of Surgery: The Biological Basis of Modern Surgical Practice. 21st ed. Elsevier; 2022. p. 897-913.
  \item Bilezikian JP, Brandi ML, Eastell R, et al. Guidelines for the Management of Asymptomatic Primary Hyperparathyroidism: Summary Statement from the Fourth International Workshop. J Clin Endocrinol Metab. 2014;99(10):3561-3569.
  \item American Association of Endocrine Surgeons. Parathyroidectomy. Available from: \url{https://www.endocrinesurgery.org/patient-education/parathyroidectomy}
\end{enumerate}

\clearpage
